\begin{lemma}[Proper prefixes of a minimum bad tree]
Let \((\varphi,T)\) be a minimum bad pair.  If \(S\) is a proper initial
segment of \(T\), then \(V(S)\) is \(\varphi\)-elementary.
\end{lemma}
\begin{proof}
Write \(q=\rho(T)\), using \noderef{TashTreePathIndex}.  The initial
segment relation in \noderef{TreeSequenceInitialSegment} identifies every
displayed vertex and exact edge of \(S\) with the corresponding initial
vertex and edge of \(T\).  Therefore every Tashkinov witness required for
an edge of \(S\) is inherited from the witness for the same edge of \(T\):
the witness occurs earlier than that edge and hence is still in the initial
segment.  Consequently \(S\) is a Tashkinov tree for the same \(k\)-triple.

We compare ranks.  If \(|E(S)|<q\), the general bound
\(\rho(S)\le |E(S)|\), which follows from
\noderef{TashTreeSuffixPath} by taking the last displayed vertex, gives
\(\rho(S)<q\).  If \(q\le |E(S)|\), the path suffix of \(T\) beginning at
\(y_q\) restricts to a path suffix of \(S\), so minimality in the
definition of \(\rho(S)\) gives \(\rho(S)\le q\).  A strict inequality
again makes the first rank coordinate smaller.  In the remaining case
\(\rho(S)=q\), while the proper-prefix hypothesis gives
\(|E(S)|<|E(T)|\).  Thus in every case
\(S\) has smaller rank than \(T\) according to
\noderef{TashTreeRankLess}.

If \(V(S)\) were not elementary, the preceding inheritance would make
\((\varphi,S)\) a bad pair of smaller rank.  This contradicts the defining
minimality clause of \noderef{TashTreeMinimumBadPair}.  Hence the proper
prefix is elementary.
\end{proof}
